\subsection{连通紧李群的自同构}

用 \(\mathrm{Aut}(\mathrm{G})\) 表示 G 的自同构所成的李群（第三章，§10，第 2 节），并用 \(\mathrm{Aut}(\mathrm{G},\mathrm{T})\) 表示 \(\mathrm{Aut}(\mathrm{G})\) 中由元素 \(u\)（满足 \(u(\mathrm{T}) = \mathrm{T}\) 者）组成的闭子群。我们已看到（§1，第 4 节，命题 4 的推论 5），\(\mathrm{Aut}(\mathrm{G})\) 的单位元连通分支是内自同构子群 \(\mathrm{Int}(\mathrm{G})\)；用 \(\mathrm{Int}_{\mathrm{G}}(\mathrm{H})\) 表示 G 的子群 H 在 \(\mathrm{Int}(\mathrm{G})\) 中的像。

设 D 是 G 关于 T 的共变图式；用 \(\operatorname{Aut}(D)\) 表示它的自同构群，并用 \(\mathrm{W}(D)\) 表示它的 Weyl 群。映射 \(u \mapsto \mathrm{D}_{*}(u)\) 是从 \(\operatorname{Aut}(G, T)\) 到 \(\operatorname{Aut}(D)\) 的同态。第 9 节命题 15 立即给出：

命题 17. 同态 \(\operatorname{Aut}(\mathrm{G},\mathrm{T})\to \operatorname{Aut}(\mathrm{D})\) 是满射，其核为 \(\operatorname{Int}_{\mathrm{G}}(\mathrm{T})\)。

注意 \(\operatorname{Aut}(G,T) \cap \operatorname{Int}(G) = \operatorname{Int}_{G}(N_{G}(T))\)，且 \(\operatorname{Int}_{G}(N_{G}(T))\) 在 \(\operatorname{Aut}(D)\) 中的像是 \(W(D)\)（第 5 节，命题 10）。于是，命题 17 给出同构

\[
\operatorname{Aut} (\mathrm{G}, \mathrm{T}) / (\operatorname{Aut} (\mathrm{G}, \mathrm{T}) \cap \operatorname{Int} (\mathrm{G})) \rightarrow \operatorname{Aut} (\mathrm{D}) / \mathrm{W} (\mathrm{D}).
\]

此外，\(\operatorname{Aut}(\mathrm{G}) = \operatorname{Int}(\mathrm{G}).\operatorname{Aut}(\mathrm{G}, \mathrm{T})\)。实际上，若 u 属于 \(\operatorname{Aut}(\mathrm{G})\)，则 \(u(\mathrm{T})\) 是 T 的一个极大环面，故与 T 共轭，从而存在 G 的内自同构 v 使得 \(u(\mathrm{T}) = v(\mathrm{T})\)，换言之，\(v^{-1}u \in \operatorname{Aut}(\mathrm{G}, \mathrm{T})\)。由此可知，\(\operatorname{Aut}(\mathrm{G})/\operatorname{Int}(\mathrm{G})\) 可等同于 \(\operatorname{Aut}(\mathrm{G}, \mathrm{T})/(\operatorname{Aut}(\mathrm{G}, \mathrm{T}) \cap \operatorname{Int}(\mathrm{G}))\)，故由前述结果，我们有正合序列

\[
1 \to \operatorname{Int} (\mathrm{G}) \to \operatorname{Aut} (\mathrm{G}) \to \operatorname{Aut} (\mathrm{D}) / \mathrm{W} (\mathrm{D}) \to 1.\tag{16}
\]

因此：

命题 18. 群 \(\operatorname{Aut}(\mathrm{G}) / \operatorname{Int}(\mathrm{G})\) 同构于 \(\operatorname{Aut}(\mathrm{D}) / \operatorname{W}(\mathrm{D})\)。

特别地，假设 G 是半单的；这时群 \(\operatorname{Aut}(D)\) 可等同于 \(\mathrm{A}(\mathrm{R}(\mathrm{G},\mathrm{T}))\)（第六章，§1，第 1 节）中由满足 \(u(\mathrm{X}(\mathrm{T})) \subset \mathrm{X}(\mathrm{T})\) 的元素 u 组成的子群，而子群 \(\mathrm{W}(\mathrm{D})\) 可等同于 \(\mathrm{W}(\mathrm{R}(\mathrm{G},\mathrm{T}))\)。

推论. 若 G 是单连通的，或 C(G) 仅由单位元素组成，则群 \(\operatorname{Aut}(\mathrm{G})/\operatorname{Int}(\mathrm{G})\) 同构于 R(G,T) 的邓金图的自同构群。

这由前述结果及第六章 §4 第 2 节命题 1 的推论得出。

我们现在要证明扩张 (16) 容有截面。

对所有 \(\alpha\in\mathrm{R}(\mathrm{G},\mathrm{T})\)，用 \(\mathrm{V}(\alpha)\) 表示 g 的满足 \(\mathrm{V}(\alpha)_{(\mathbf{C})}=\mathfrak{g}^{\alpha}+\mathfrak{g}^{-\alpha}\) 的 2 维向量子空间；并用 K 表示相伴于 g 的 Killing 型的二次型。

定义 3. (G,T) 的一个标架是一个偶 \((\mathrm{B},(U_{\alpha})_{\alpha \in \mathrm{B}})\)，其中 B 是 \(\mathrm{R}(\mathrm{G},\mathrm{T})\) 的一个基（第六章，§1，第 5 节，定义 2），且对所有 \(\alpha \in \mathrm{B}\)，\(U_{\alpha}\) 是 \(\mathrm{V}(\alpha)\) 中满足 \(\mathrm{K}(U_{\alpha}) = -1\) 的元素。

G 的一个标架是指 G 的一个极大环面 T 连同 \((G, T)\) 的一个标架。

引理 3. 设 \(B_{0}\) 是 \(\mathrm{R}(\mathrm{G},\mathrm{T})\) 的一个基。群 \(\operatorname{Int}_{\mathrm{G}}(\mathrm{T})\) 单传递地作用于\((\mathrm{G},\mathrm{T})\) 的形如 \((\mathrm{B}_{0},(U_{\alpha})_{\alpha\in\mathrm{B}_{0}})\) 的标架所成的集合上。

对所有 \(\alpha\in B_{0}\)，用 \(\mathrm{K}(\alpha)\) 表示二次型 K 在 \(\mathrm{V}(\alpha)\) 上的限制；T 在 \(\mathrm{V}(\alpha)\) 上的作用定义一个态射 \(\iota_{\alpha}:T\to\mathbf{SO}(\mathrm{K}(\alpha))\)。我们在第 4 节已看到，\(\mathbf{SO}(\mathrm{K}(\alpha))\) 可等同于 U，使得 \(\iota_{\alpha}\) 恰好等同于根 \(\alpha\)。由于 \(B_{0}\) 是 R 的基，它是由根生成的 Z-模 \(\mathrm{Q}(\mathrm{R})\) 的基，因而是 \(\mathrm{X}(\mathrm{T}/\mathrm{C}(\mathrm{G}))\)（\(\mathrm{X}(\mathrm{T})\) 的一个子模）的基。由此可知，诸态射 \(\iota_{\alpha}\) 的乘积诱导出从 \(\mathrm{T}/\mathrm{C}(\mathrm{G})\) 到诸群 \(\mathbf{SO}(\mathrm{K}(\alpha))\) 的乘积的同构。而后一群单传递地作用于第一分量为 \(B_{0}\) 的 (G,T) 的标架所成的集合上。

命题 19. 群 \(\operatorname{Int}(\mathbf{G})\) 单传递地作用于 \(\mathbf{G}\) 的标架所成的集合上。

设 \(e = (\mathrm{T}, \mathrm{B}, (U_{\alpha}))\) 与 \(e' = (\mathrm{T}', \mathrm{B}', (U_{\alpha}')\) 是 \(G\) 的两个标架。存在元素 \(g\) 属于 \(G\)，使得 \((\operatorname{Int} g)(\mathrm{T}) = \mathrm{T}'\)，且这些元素组成模 \(N_G(T)\) 的单个陪集。于是可以假设 \(T = T'\)，而我们必须证明存在 \(\operatorname{Int}_G(N_G(T))\) 的唯一元素把 \(e\) 变到 \(e'\)。由第六章 §1 第 5 节注 4，存在唯一的 \(w\) 属于 \(W(R)\)，使得 \(w(B) = B'\)。由于 \(W(R)\) 可等同于 \(N_G(T)/T\)，故存在 \(n \in N_G(T)\) 使得 \(w = \operatorname{Int} n\)，且 \(n\) 模 \(T\) 唯一确定。于是可以假设 \(B = B'\)，而我们必须证明存在 \(\operatorname{Int}_{\mathrm{G}}(\mathrm{T})\) 的唯一元素把 e 变到 \(e'\)，这正是引理 3。

推论. 设 e 是 (G,T) 的一个标架，并设 E 是使 e 稳定的 G 的自同构群。则 \(\operatorname{Aut}(G)\) 是 E 与 \(\operatorname{Int}(G)\) 的半直积，而 \(\operatorname{Aut}(G,T)\) 是 E 与 \(\operatorname{Int}(G)\cap\operatorname{Aut}(G,T)=\operatorname{Int}_{G}(\mathrm{N}_{\mathrm{G}}(\mathrm{T}))\) 的半直积。

实际上，\(\operatorname{Aut}(G)\) 的每个元素把 e 变为 G 的一个标架。由命题 19，\(\operatorname{Aut}(G)\) 模 \(\operatorname{Int}(G)\) 的每个陪集与 E 恰交于一点，由此得第一个断言。第二个断言同理可证。

注. 设 G 与 \(G'\) 是两个连通紧李群，并设 \(e = (\mathrm{T}, \mathrm{B}, (U_{\alpha}))\) 与 \(e' = (\mathrm{T}', \mathrm{B}', (U_{\alpha}')\) 分别是 G 与 \(G'\) 的标架。设 X 是从 G 到 \(G'\) 且把 e 变到 \(e'\) 的同构全体所成的集合。映射 \(f \mapsto \mathrm{D}^{*}(f)\)（相应地 \(\mathrm{D}_{*}(f)\)）是从 X 到这样的同构全体所成集合的双射：这些同构是从 \(\mathrm{D}^{*}(\mathrm{G}', \mathrm{T}')\) 到 \(\mathrm{D}^{*}(\mathrm{G}, \mathrm{T})\)（相应地，从 \(\mathrm{D}_{*}(\mathrm{G}, \mathrm{T})\) 到 \(\mathrm{D}_{*}(\mathrm{G}', \mathrm{T}')\)）的同构，且把 \(B'\) 变到 B（相应地，把 B 变到 \(B'\)）。实际上，这由命题 15 与引理 3 立即得出。

